\subsection{The actual global presentation and the coupled local packets}

\begin{statement}{E.73}{Proved}{The complete relation set for global local-packet gluing}
Put $K=\mathbb Q(\zeta_3)$, $S=\{(1-\zeta_3)\}$ and
\[
 \mathcal G=G_K^{(3)},\qquad
 P=G_{K,S}^{(3)}=\langle\tau,\tau'\rangle,
 \qquad L=\ker(\mathcal G\to P).
 \tag{E73.1}
\]
Use the actual splitting $s:P\to\mathcal G$ of E.39.
For each finite prime $\mathfrak q\notin S$, choose an actual
prolongation, a positive tame generator $t_{\mathfrak q}$ and
an arithmetic Frobenius $F_{\mathfrak q}$ in its decomposition
group. Write $q=N\mathfrak q$. Let $\mathcal F$ be the free
pro-three group on $\tau,\tau'$ and the pointed convergent
family $(T_{\mathfrak q})_{\mathfrak q\notin S}$, with
$T_{\mathfrak q}\to1$. The continuous map
\[
 \pi:\mathcal F\longrightarrow\mathcal G,
 \qquad \tau\mapsto s(\tau),\quad
 \tau'\mapsto s(\tau'),\quad T_{\mathfrak q}\mapsto t_{\mathfrak q}
 \tag{E73.2}
\]
is a minimal epimorphism. Choose actual preimages
$Y_{\mathfrak q}\in\mathcal F$ with
$\pi(Y_{\mathfrak q})=F_{\mathfrak q}$. Then its entire kernel is
\[
 \ker\pi=
 \overline{\langle\!\langle
   Y_{\mathfrak q}T_{\mathfrak q}Y_{\mathfrak q}^{-1}
        T_{\mathfrak q}^{-q}:\mathfrak q\notin S
 \rangle\!\rangle}.
 \tag{E73.3}
\]
This is a presentation with actual pro-three words
$Y_{\mathfrak q}$, not prescribed independent Frobenius symbols.
Its cohomology and relation module have the exact topologies
\[
 \begin{split}
 H^2_{\rm cts}(\mathcal G,\mathbb F_3)
   &\xrightarrow[\mathrm{restriction}]{\ \sim\ }
      \bigoplus_{\mathfrak q\notin S}
              H^2_{\rm cts}(D_{\mathfrak q},\mathbb F_3)
      \simeq\bigoplus_{\mathfrak q\notin S}\mathbb F_3,\\
 R/\overline{R^3[R,\mathcal F]}
   &\simeq\prod_{\mathfrak q\notin S}\mathbb F_3,
       \qquad R=\ker\pi.
 \end{split}
 \tag{E73.4}
\]
The relators displayed in (E73.3) are the coordinate vectors
of the compact product in (E73.4). No product replaces
the continuous-cohomology direct sum.

Let $V$ be one of the actual geometric pro-three groups
$\Pi$, $\Pi/\Pi''$, or $\Pi/\gamma_3 I_Q$, with its given
based action $\rho:P\to\operatorname{Aut}(V)$. Define
\[
 c_0(\mathcal P_K\setminus S,V)
   =\{(n_{\mathfrak q}): n_{\mathfrak q}\to1\}.
 \tag{E73.5}
\]
For $A,A'\in V$ and such a family, the universal property
gives the unique homomorphism
\[
 \widetilde a:\mathcal F\longrightarrow V\rtimes\mathcal G,
 \quad
 \tau\mapsto(A,s\tau),\quad
 \tau'\mapsto(A',s\tau'),\quad
 T_{\mathfrak q}\mapsto(n_{\mathfrak q},t_{\mathfrak q}).
 \tag{E73.6}
\]
Define $v_{\mathfrak q}$ by the actual word evaluation
$\widetilde a(Y_{\mathfrak q})=(v_{\mathfrak q},F_{\mathfrak q})$.
These data give a global section exactly when, for every
$\mathfrak q\notin S$,
\[
 v_{\mathfrak q}\,
     \rho(\sigma_{\mathfrak q})(n_{\mathfrak q})
         v_{\mathfrak q}^{-1}=n_{\mathfrak q}^{q},
 \qquad \sigma_{\mathfrak q}=\operatorname{pr}(F_{\mathfrak q}).
 \tag{E73.7}
\]
Thus the actual local equations are sufficient as well
as necessary; there is no further relation after all
the evaluated words in (E73.7) have been retained.
Their local solution packets are exactly E.61, E.66,
E.68 and E.71. In particular a ramified full or
metabelian packet has, for one of the nine actual
cusps, the form
\[
 \begin{split}
 n_{\mathfrak q}&=g_{\mathfrak q}d_i^a g_{\mathfrak q}^{-1},\\
 v_{\mathfrak q}&=
    g_{\mathfrak q}d_i^b z_{i,\mathfrak q}
       \rho(\sigma_{\mathfrak q})(g_{\mathfrak q})^{-1},
       \qquad a\ne0,
 \end{split}
 \tag{E73.8}
\]
where $F_{i,\mathfrak q}=z_{i,\mathfrak q}F_{\mathfrak q}$
is the actual tangential lift. For the marked quotient,
the $Q$ packet additionally has the central coordinate
of E.67 and E.70. The unramified packet has
$n_{\mathfrak q}=1$ and arbitrary $v_{\mathfrak q}$.

Let $\mathcal X_V$ denote the subset of
$V^2\times c_0(\mathcal P_K\setminus S,V)$ defined by
(E73.7). Restriction to the displayed generators gives
the homeomorphism
\[
 \operatorname{Sect}_K(V)/V\simeq\mathcal X_V/V,
 \tag{E73.9}
\]
with its full conjugation action
\[
 (A,A',n_{\mathfrak q})\longmapsto
 \bigl(wA\rho(\tau)(w)^{-1},\,
       wA'\rho(\tau')(w)^{-1},\,
       wn_{\mathfrak q}w^{-1}\bigr).
 \tag{E73.10}
\]
The topology on $c_0$ is uniform congruence at every
coordinate, equivalently the inverse limit of the
discrete finite-support configuration spaces for
the finite quotients of $V$. The quotient maps from
$\Pi$ to the two approximations induce exactly the
componentwise maps in (E73.9), with their word evaluations
and local packet maps retained.

For sections over $\mathbb Q$, there is an equally exact
descent condition. Fix the actual complex conjugation $c$.
Choose words $C_{\mathfrak q}\in\mathcal F$ lifting
$ct_{\mathfrak q}c$, tending to one, and take
$C_\tau=\tau'$, $C_{\tau'}=\tau$. Write
$\widetilde a(C_{\mathfrak q})
=(a_{\mathfrak q}^{c},ct_{\mathfrak q}c)$.
After conjugating the value at $c$ to one, the precise
additional equations and remaining conjugations are
\[
 A'=c(A),\qquad a_{\mathfrak q}^{c}=c(n_{\mathfrak q})
       \quad(\mathfrak q\notin S),\qquad w\in V^c.
 \tag{E73.11}
\]
No chosen inertia prolongation is presumed fixed by $c$.

Independent local packets do not give a Cartesian product
of global parameters. In the actual cubic Kummer quotient
of (E.39), the prescribed values on $s(\tau),s(\tau')$
and inertia at two are $(2,1,1)\in\mathbb F_3^3$,
with zero values on all other inertia generators.
They force the value at $(-2+3\zeta)$ over nineteen to
be one. Its quotient-section Frobenius value is two.
Both values define valid unramified local packets,
and only the former extends with those global generator
values. This is an actual finite obstruction to omitting
the word evaluation in (E73.7), verified by the retained
program. The theorem leaves the evaluation of all
$Y_{\mathfrak q}$, and the solution of the resulting
coupled global equations, as the stated arithmetic
input; it does not identify that unevaluated system
with an independent product of the local classes.

\emph{Manuscript: \S51.16, folios733v--736v, scans172--178,
formulas(371)--(379); Part E(c),(d) and follow-up(4).}
\end{statement}

\begin{proof}
\emph{The actual inertia bundle and its support.}
E.39, using Wingberg's Theorem2.4, gives the continuous
free product of the inertia groups over the primes of
$K_S(3)$. The preceding definition and Lemma1.1 of that
same original give its cohomology with the continuous
direct-sum topology
\cite[pp.~4--5, Lemma1.1; pp.~13--14, Theorem2.4]{E_WingbergRiemann}.
Here is the specialization, including all support maps.
For a chosen prime over $\mathfrak q$ let
$\Delta_{\mathfrak q}\subset P$ be its decomposition
subgroup. It is $\mathbb Z_3$: the unramified cyclotomic
character sends its Frobenius to $q\ne1$ in
$1+3\mathbb Z_3$. Its actual local group in $\mathcal G$
is the full tame pro-three group
\[
 D_{\mathfrak q}=\langle t_{\mathfrak q}\rangle
                    \rtimes\langle F_{\mathfrak q}\rangle,
 \qquad F_{\mathfrak q}t_{\mathfrak q}F_{\mathfrak q}^{-1}
          =t_{\mathfrak q}^{q}.
 \tag{E73.12}
\]
Indeed its inertia injection is supplied by the actual
free product; its unramified quotient injects because
the displayed cyclotomic character has infinite order.

Put $M=H^1_{\rm cts}(L,\mathbb F_3)$, a discrete
$P$-module. Restriction to the inertia factors gives
\[
 M=\bigoplus_{\mathfrak q\notin S}
       \operatorname{Coind}_{\Delta_{\mathfrak q}}^P\mathbb F_3.
 \tag{E73.13}
\]
The summand is the locally constant functions on
$P/\Delta_{\mathfrak q}$, with the contragredient
permutation action. The mod-three orientation is
trivial because $\mu_3\subset K$. A continuous
homomorphism from $L$ to the finite group $\mathbb F_3$
is nonzero over only finitely many base primes and is
locally constant on each of their prolongation spaces.
Conversely those data give the continuous homomorphism
by the free-product universal property. This proves
(E73.13), including the direct-sum assertion. It also
identifies its $P$-invariants with the direct sum of
the constant functions on those spaces.

Both $P$ and $L$ have cohomological dimension one.
The split extension $L\rtimes P$ consequently gives
\[
 H^2(\mathcal G,\mathbb F_3)=H^1(P,M).
 \tag{E73.14}
\]
We describe this edge map to retain its naturality.
A central extension of $\mathcal G$ by $\mathbb F_3$
splits continuously on $L$ and on $s(P)$. Comparing
the two lifted conjugation actions gives a one-cocycle
$z:P\to\operatorname{Hom}_{\rm cts}(L,\mathbb F_3)$.
Changing the splitting on $L$ changes $z$ by a
coboundary. Conversely such a $z$ gives the explicit
two-cocycle, in the coordinates $(\ell,p)\in L\rtimes P$,
\[
 \omega_z((\ell,p),(\ell',p'))
               =z(p)(p\ell'p^{-1}).
 \tag{E73.15}
\]
The cocycle equation for $z$ and additivity on $L$
verify the two-cocycle equation directly. The two
constructions are inverse. They are continuous:
the conjugation discrepancy is a continuous function
on the compact product $P\times L$ with finite range,
so it is locally constant uniformly in its $L$
coordinate. This proves (E73.14) as well as its
compatibility with actual subgroup restrictions.

Every continuous $P$-cocycle with values in (E73.13)
has finite support in the base primes. For example its
values on the two free generators lie in a finite
sum of summands, which is $P$-stable and hence
contains all its values. The same assertion holds
for coboundaries. Thus $H^1(P,-)$ commutes with this
particular direct sum. The Shapiro map on each
summand is restriction to $\Delta_{\mathfrak q}$
followed by evaluation at its base coset.
For precision, a continuous transversal of
$P\to P/\Delta_{\mathfrak q}$ is constructed by
compatible choices in the finite quotients of $P$.
Each prescribed finite representative lifts to the
next level after multiplying by an element of
$\Delta_{\mathfrak q}$. If $j(x)$ is this transversal,
with $j(\Delta)=1$, the inverse for a character
$\psi:\Delta\to\mathbb F_3$ is
\[
 z_\psi(p)(x)
   =\psi\bigl(j(x)^{-1}p\,j(p^{-1}x)\bigr).
 \tag{E73.16}
\]
The argument lies in $\Delta$. Multiplication verifies
the cocycle equation. For any cocycle $z$, setting
$m(x)=z(j(x))(x)$ gives
$z=z_\psi+m-pm$, proving both directions. Compactness
of $P/\Delta$ gives the uniform local constancy needed
for a discrete coinduced module.

The local analogue of (E73.14), applied to (E73.12),
identifies $H^2(D_{\mathfrak q},\mathbb F_3)$ with
$H^1(\Delta_{\mathfrak q},\mathbb F_3)=\mathbb F_3$.
In the central extension described above, lifting
$F_{\mathfrak q}=h_{\mathfrak q}s(\sigma_{\mathfrak q})$
as the product of its two chosen lifts shows that
its local relation error is
\[
                  z(\sigma_{\mathfrak q})(t_{\mathfrak q}).
 \tag{E73.17}
\]
Inner conjugation by $h_{\mathfrak q}\in L$ disappears
on the character of $L$, and the orientation factor
$q$ is one in $\mathbb F_3$. Thus (E73.17) is precisely
the Shapiro evaluation, with the positive sign for
the relator in (E73.3). This proves the restriction
isomorphism in (E73.4), with finite support, directly.

\emph{Minimal generators and the complete relation set.}
The analogous degree-one restriction sequence, split
by $s$, gives the exact evaluation isomorphism
\[
 H^1(\mathcal G,\mathbb F_3)
   \longrightarrow\mathbb F_3^2\oplus
                  \bigoplus_{\mathfrak q\notin S}\mathbb F_3,
 \quad \chi\longmapsto
    (\chi(s\tau),\chi(s\tau'),\chi(t_{\mathfrak q})).
 \tag{E73.18}
\]
Every finite extension of $K$ ramifies at finitely
many primes, so the chosen $t_{\mathfrak q}$ tend
to one in $\mathcal G$, independently of the
chosen prolongations. This proves continuity of
(E73.2). If its compact image were proper, a finite
three-group quotient would contain that image in a
proper subgroup and then in a subgroup of index
three. The resulting nonzero character would
vanish on every displayed generator, contradicting
(E73.18). Thus $\pi$ is onto, and (E73.18) proves
that it induces an isomorphism on Frattini quotients.

For a central extension of $\mathcal G$ by
$\mathbb F_3$, lift the free group $\mathcal F$
to that extension. These lifts respect the convergent
basis: an open subgroup of the extension avoiding its
finite kernel maps isomorphically onto an open subgroup
of $\mathcal G$, and contains the unique small lifts
of all sufficiently late inertia generators.
Restriction of the lift to $R$
is an $\mathcal F$-invariant character. Different
lifts change it by the restriction of a character
of $\mathcal F$, which is zero on $R$ by minimality.
Conversely pushout by such an invariant character
constructs the central extension. Consequently
\[
 H^2(\mathcal G,\mathbb F_3)
      =\operatorname{Hom}_{\rm cts}
          (R/\overline{R^3[R,\mathcal F]},\mathbb F_3).
 \tag{E73.19}
\]
This is the actual relation-module map: its value
on a relator is the central coordinate obtained by
evaluating that word in any lift of the generators.
Pulling back along the two-generator local free
group, with its generators sent to
$T_{\mathfrak q},Y_{\mathfrak q}$, identifies this
evaluation with the local coordinate (E73.17).
The tame group (E73.12) itself has the one-relator
presentation in (E73.3): its relation makes the
closed cyclic inertia normal, and collection
gives $t^aF^b$, whose two coordinates are detected
in $\mathbb Z_3\rtimes_q\mathbb Z_3$.

It follows from (E73.4) and (E73.19) that the
displayed relators span the compact relation
module densely, and that their evaluations are
the coordinate basis. They tend to one in
$\mathcal F$: for every open normal subgroup,
eventually $T_{\mathfrak q}$ belongs to it,
as do its conjugate and every integral power.
This argument needs no convergence of the
$Y_{\mathfrak q}$ themselves.

Let $R_0$ be their closed normal closure. Were
$R_0\ne R$, the nontrivial normal subgroup
$H=R/R_0$ of $\mathcal F/R_0$ would have a
nontrivial finite three-group image $J$ in a
finite quotient. The finite module $J/\Phi(J)$
has nonzero coinvariants under that finite
three-group: the augmentation ideal in its
group algebra over $\mathbb F_3$ is nilpotent.
It therefore gives a nonzero invariant character
of $H$, hence of $R$, vanishing on $R_0$.
This contradicts dense spanning in (E73.19).
Thus $R_0=R$, proving (E73.3).

The same proof applies with any finite set of
additional primes. Deleting the other inertia
generators sends $Y_{\mathfrak q}$ to actual
Frobenius preimages in that finite-ramification
presentation. All omitted relators become one.
The inverse limit of those actual quotient maps
recovers (E73.2)--(E73.3). This also proves the
stated product topology on the relation module.

\emph{Sufficiency and the nonlinear arithmetic correction.}
The actual action on each $V$ factors through
$P$ by E.34, so $L$ acts trivially. Its semidirect
product with $\mathcal G$ is pro-three. More explicitly,
a cocycle on the full arithmetic group $G_K$ defines
a homomorphism to $V\rtimes\rho(G_K)$, whose target
is pro-three. It therefore factors through
$\mathcal G$, and conversely every constructed
cocycle inflates to $G_K$. These are inverse maps
on the actual based sections and their conjugacies.
The
convergence condition (E73.5) is exactly what
makes (E73.6) continuous on the pointed generating
space. Evaluation of a relator in (E73.3) has
geometric component
\[
 v_{\mathfrak q}\rho(\sigma_{\mathfrak q})(n_{\mathfrak q})
       v_{\mathfrak q}^{-1}n_{\mathfrak q}^{-q}.
 \tag{E73.20}
\]
Thus all (E73.7) hold exactly when the homomorphism
(E73.6) kills the entire kernel, proving sufficiency.
Its descended homomorphism has second projection
the identity on $\mathcal G$, so it is the desired
section. Restriction is the inverse construction.

Let $u:P\to V$ be the free cocycle specified by
$A,A'$. In $\mathcal F$ the word
$Y_{\mathfrak q}s_{\mathcal F}(\sigma_{\mathfrak q})^{-1}$
lies in the kernel of $\mathcal F\to P$ and maps
to $h_{\mathfrak q}$ of E.39. On a descended solution,
its evaluation is $\lambda(h_{\mathfrak q})$. Hence
the actual equality is
\[
       v_{\mathfrak q}
             =\lambda(h_{\mathfrak q})u(\sigma_{\mathfrak q}).
 \tag{E73.21}
\]
Before descent, (E73.6) evaluates the displayed
free word itself. This specifies the correction
without presupposing a homomorphism $\lambda$
or its equivariance. The entire global relation
problem is consequently the coupled equations
(E73.7), not independently chosen local Frobenius
coordinates. E.61, E.66 and E.71 applied to the
actual pair $(n_{\mathfrak q},v_{\mathfrak q})$
give exactly the packets in the statement.

Choose cofinal characteristic open subgroups
$V_m$ of $V$, so the action preserves them.
Two cocycles are uniformly equal modulo $V_m$
exactly when their values agree on all the
displayed generators. The implication from
generators follows by their images in the
finite semidirect product using the actual
finite arithmetic action on $V/V_m$.
Thus restriction gives the claimed homeomorphism
before conjugacy. Word evaluation is continuous
for these topologies, so $\mathcal X_V$ is closed
in $V^2\times c_0$. Formula (E73.10) is the
literal section-conjugation formula; taking
its quotient proves (E73.9). The compact group
$V$ acts continuously on a Hausdorff space, so
this orbit space is Hausdorff. The actual
geometric quotient maps commute with every
generator and word evaluation, proving their
componentwise description.

\emph{Rational descent with all conjugators.}
The actual arithmetic quotient over $\mathbb Q$
is $\mathcal G\rtimes\langle c\rangle$. Every
section can be conjugated so that its value at
$c$ is one, by coprime complement conjugacy in
the extension by the pro-three group $V$.
The remaining conjugations are exactly $V^c$.
The elements $ct_{\mathfrak q}c$ tend to one.
Since the epimorphism $\pi$ is open, compatible
choices in successively smaller open subgroups
of $\mathcal F$ give preimages
$C_{\mathfrak q}\to1$. For the two free
generators, the actual splitting is equivariant,
so the indicated preimages are exactly
$\tau'$ and $\tau$.

For a descended section, (E73.11) says that the
two continuous homomorphisms obtained by conjugating
the arithmetic source, or by conjugating both
factors of the semidirect target, agree on every
generator. They therefore agree everywhere.
This is precisely the extension condition over
$\mathbb Q$. Conversely every normalized rational
section has these equations. Changing any
preimage $C_{\mathfrak q}$ or $Y_{\mathfrak q}$
by an element of $R$ changes none of the descended
evaluations. The coordinate description is
therefore independent of these choices by the
identity map on its generator values; the actual
prolongations and words remain part of its typed
construction.

\emph{An actual finite failure of independent Frobenius packets.}
We use the explicit Kummer extension
\[
 K(\sqrt[3]{2},\sqrt[3]{\zeta},\sqrt[3]{1-\zeta}).
\]
The three Kummer classes are independent: the
valuations at two and at $1-\zeta$ separate the
first and third, while $\zeta$ is not a cube in
$K$. Its Galois group is consequently
$H=\mathbb F_3^3$, with coordinates
$(\kappa_2,\kappa_\zeta,\kappa_{1-\zeta})$.
The actual Frobenii and the inertia generator
give
\[
 \begin{array}{c|ccc}
  &\kappa_2&\kappa_\zeta&\kappa_{1-\zeta}\\\hline
  s\tau&2&2&0\\
  s\tau'&1&2&2\\
  t_{(2)}&1&0&0\\
  F_{(-2+3\zeta)}&1&0&2.
 \end{array}
 \tag{E73.22}
\]
The first two rows are the exact E.39 choice
of the Frobenii above seven. The last is
obtained by reduction with $\zeta=7$ in
$\mathbb F_{19}$. The projection of that
last row to $s(P)$ is
$2s\tau+s\tau'=(2,0,2)$, and their difference
is $2t_{(2)}$. The unique character taking
the prescribed values $(2,1,1)$ on the first
three rows is $\kappa_2$, so its value on the
last row is one, not two.

This finite obstruction is detected by an actual
geometric quotient. With the original tangent
at $O=C_0$, the unit $z/x$ on the affine cubic
induces the based map $\Pi\to\mathbb Z_3(1)$;
its order at $O$ is one. Reduction modulo three
is equivariant over $K$ and factors through
both approximations. The cusp cocycle
$g\mapsto d_O^{\kappa_2(g)}$ realizes the first
three rows and the value one at nineteen.
Keeping its generator data and replacing only
that unramified local value by two remains a
valid local packet but violates (E73.22).
Conjugacy cannot change this finite character
because its coefficient action over $K$ is
trivial. The program recomputes all residues,
enumerates all27 characters of $H$, and verifies
the forced value and the actual correction.
It also checks the explicit Shapiro cochains
and the positive local-relator transgression
in finite groups. This proves every asserted
comparison while retaining the unevaluated
higher Frobenius words in (E73.21).
\end{proof}
